\documentclass[12pt]{article}
\usepackage{natbib}

%opening
\title{GenAI is the Enemy of Education}
\author{Tess O'Brien}

\begin{document}

\maketitle



\section{Introduction: The current moment}
Since the release of ChatGPT in 2022 \citep{chatgptRelease2022}, generative artificial intelligence (GenAI) in the form of large language models (LLMs) products have exploded into education across multiple fields, particularly the sciences. The effects of incorporating GenAI into teaching have been sold to teachers and students as revolutionary and positive. We are told that our students will be `left behind' if they do not learn to use GenAI tools, that we can (and should) use GenAI to replace as many aspects of the teaching process as possible in order to maximise the number of students we can teach and the amount of material we can give them.

I reject these claims. I believe that we are seeing GenAI negatively impact both the learning of our students and the teaching process. This article serves as a summary of the current scientific evidence for negative effects of GenAI on skills in the workforce, learning, and teaching, as well as some more hypothetical arguments for the likely downstream impacts. 

Many of the articles referenced here are pre-prints, which means that they have not been subject to the peer review process of typical publication. I assume that they are published under good faith and reasonable practices, and will critique the methods as I see necessary.

\section{The goal is to think}
- Outsourcing the hard work of thinking to an LLM means students fail to develop that ability

- Summarising is a key pedagogical tool.

- Developing expertise has no shortcuts.


\section{Entry-level problems matter}
- Replacing grunt-work with LLMs blows a hole in the training pipeline for programmers and data analysts.

\section{Unreliability matters}
- Probabilistic models such as LLMs are particularly bad for processes that require consistency, such as marking.

\section{Teaching is a social process}
- We lose so much by introducing a barrier between ourselves and our students.

- The student-teacher relationship is the biggest value-add of higher education.

- A theory of mind is necessary to teach.

- Marking is a conversation with the student and should not be outsourced to machines.

\newpage
\section{Acknowledgements}
In no particular order, here is a list of names and usernames for people who have contributed sources or listened to me talk about this project: AnastasiaGals, deltaceti, urizenxvii, \_cnorm.

\newpage
\bibliographystyle{apalike}
\bibliography{article}
\end{document}
